% GENERATED FILE. Edit canonical lessons, then run npm run generate:books.
% canonical-source-hash: fnv1a64:440ed1d7acbe03ee
% canonical-lessons: HI-C200-tvaca, HI-C200-rajya, HI-C200-rajdhani, HI-C200-sima, HI-C200-dvip

\chapter{Skin, A state, A capital city, A border, An island}
\label{ch:hi-skin-a-state-a-capital-city-a-border-an-island}
\hlchaptermodality{200}

\begin{chapteropening}
\textbf{By the end of this chapter:} \emph{I can say skin, a state, a capital city, a border and an island.}

Complete the last lesson of chapter 200: I can say skin, a state, a capital city, a border and an island.
\end{chapteropening}

\section[\texorpdfstring{\hi{त्वचा} (tvacā)}{त्वचा (tvacā)}]{\hi{त्वचा} (tvacā) — skin}

\label{lesson:HI-C200-tvaca}



\begin{quote}
Before the new one: say the Hindi for the heart, then the Hindi for a bone.
\end{quote}

\subsection*{You'll want to know: \hi{त्वचा}}
\textbf{\hi{त्वचा}} — \emph{tvacā} — \textquotedblleft{}skin\textquotedblright{}.

\begin{grammarlens}[title={What you've built}]
One more word for this chapter.
\end{grammarlens}

\subsection*{Your turn}
\begin{itemize}
\raggedright
  \item \emph{Say it:} \emph{tvacā}
  \item \emph{Say it:} \emph{tvacā}, once more
  \item \emph{Say it:} say \emph{tvacā}
  \item \emph{Recall:} say the Hindi for the heart, then the Hindi for a bone, then say \emph{tvacā} again
\end{itemize}

\begin{tcolorbox}[breakable,colback=teal!4,colframe=teal!35!black,title={Before you move on}]
What does \hi{त्वचा} mean? (\textquotedblleft{}Skin\textquotedblright{}.) Say it once more.
\end{tcolorbox}
\section[\texorpdfstring{\hi{राज्य} (rājya)}{राज्य (rājya)}]{\hi{राज्य} (rājya) — a state}

\label{lesson:HI-C200-rajya}



\begin{quote}
Before the new one: say the Hindi for a bone, then the Hindi for skin.
\end{quote}

\subsection*{You'll want to know: \hi{राज्य}}
\textbf{\hi{राज्य}} — \emph{rājya} — \textquotedblleft{}a state\textquotedblright{}.

\begin{grammarlens}[title={What you've built}]
One more word for this chapter.
\end{grammarlens}

\subsection*{Your turn}
\begin{itemize}
\raggedright
  \item \emph{Say it:} \emph{rājya}
  \item \emph{Say it:} \emph{rājya}, once more
  \item \emph{Say it:} say \emph{rājya}
  \item \emph{Recall:} say the Hindi for a bone, then the Hindi for skin, then say \emph{rājya} again
\end{itemize}

\begin{tcolorbox}[breakable,colback=teal!4,colframe=teal!35!black,title={Before you move on}]
What does \hi{राज्य} mean? (\textquotedblleft{}A state\textquotedblright{}.) Say it once more.
\end{tcolorbox}
\section[\texorpdfstring{\hi{राजधानी} (rājdhānī)}{राजधानी (rājdhānī)}]{\hi{राजधानी} (rājdhānī) — a capital city}

\label{lesson:HI-C200-rajdhani}



\begin{quote}
Before the new one: say the Hindi for skin, then the Hindi for a state.
\end{quote}

\subsection*{You'll want to know: \hi{राजधानी}}
\textbf{\hi{राजधानी}} — \emph{rājdhānī} — \textquotedblleft{}a capital city\textquotedblright{}.

\begin{grammarlens}[title={What you've built}]
One more word for this chapter.
\end{grammarlens}

\subsection*{Your turn}
\begin{itemize}
\raggedright
  \item \emph{Say it:} \emph{rājdhānī}
  \item \emph{Say it:} \emph{rājdhānī}, once more
  \item \emph{Say it:} say \emph{rājdhānī}
  \item \emph{Recall:} say the Hindi for skin, then the Hindi for a state, then say \emph{rājdhānī} again
\end{itemize}

\begin{tcolorbox}[breakable,colback=teal!4,colframe=teal!35!black,title={Before you move on}]
What does \hi{राजधानी} mean? (\textquotedblleft{}A capital city\textquotedblright{}.) Say it once more.
\end{tcolorbox}
\section[\texorpdfstring{\hi{सीमा} (sīmā)}{सीमा (sīmā)}]{\hi{सीमा} (sīmā) — a border, a limit}

\label{lesson:HI-C200-sima}



\begin{quote}
Before the new one: say the Hindi for a state, then the Hindi for a capital city.
\end{quote}

\subsection*{You'll want to know: \hi{सीमा}}
\textbf{\hi{सीमा}} — \emph{sīmā} — \textquotedblleft{}a border, a limit\textquotedblright{}.

\begin{grammarlens}[title={What you've built}]
One more word for this chapter.
\end{grammarlens}

\subsection*{Your turn}
\begin{itemize}
\raggedright
  \item \emph{Say it:} \emph{sīmā}
  \item \emph{Say it:} \emph{sīmā}, once more
  \item \emph{Say it:} say \emph{sīmā}
  \item \emph{Recall:} say the Hindi for a state, then the Hindi for a capital city, then say \emph{sīmā} again
\end{itemize}

\begin{tcolorbox}[breakable,colback=teal!4,colframe=teal!35!black,title={Before you move on}]
What does \hi{सीमा} mean? (\textquotedblleft{}A border, a limit\textquotedblright{}.) Say it once more.
\end{tcolorbox}
\section[\texorpdfstring{\hi{द्वीप} (dvīp)}{द्वीप (dvīp)}]{\hi{द्वीप} (dvīp) — an island}

\label{lesson:HI-C200-dvip}



\begin{quote}
Before the new one: say the Hindi for a capital city, then the Hindi for a border.
\end{quote}

\subsection*{You'll want to know: \hi{द्वीप}}
\textbf{\hi{द्वीप}} — \emph{dvīp} — \textquotedblleft{}an island\textquotedblright{}.

\begin{grammarlens}[title={What you've built}]
One more word for this chapter.
\end{grammarlens}

\subsection*{Your turn}
\begin{itemize}
\raggedright
  \item \emph{Say it:} \emph{dvīp}
  \item \emph{Say it:} \emph{dvīp}, once more
  \item \emph{Say it:} say \emph{dvīp}
  \item \emph{Recall:} say the Hindi for a capital city, then the Hindi for a border, then say \emph{dvīp} again
\end{itemize}

\begin{tcolorbox}[breakable,colback=teal!4,colframe=teal!35!black,title={Before you move on}]
What does \hi{द्वीप} mean? (\textquotedblleft{}An island\textquotedblright{}.) Say it once more.
\end{tcolorbox}
